\paragraph{随机方程具体表示什么。}
$X_t\in\R^2$表示位置，$m_t\in\{1,2,3\}$表示未观测的模型分量。
$\Theta_m\in\R^{4\times2}$将截距、两个位置坐标与太阳代理映射为东西／南北漂移
$b_m=\Theta_m^\top\psi$，即条件平均位移率，单位为米／保存秒。
$W_t$是二维标准Wiener过程，时长$\Delta$内的独立增量协方差为$\Delta I_2$；
$L_m$将其映射为位置噪声。因此$Q_m=L_mL_m^\top$是扩散协方差，
单位为平方米／保存秒，$L_m$单位为米／保存秒的平方根。
确定性部分表示拟合的平均趋势，随机部分表示剩余变化；没有另行观测或拟合意图、速度过程。

冻结漂移与扩散时，局部增量可用下式解释：
\begin{equation}\label{eq:local-increment}
 X_{t+\Delta}-X_t\mid(X_t,m_t=m)
 \ \approx\ \mathcal N\!\left(b_m\Delta,Q_m\Delta\right).
\end{equation}
平均位移随$\Delta$增长，随机标准差随$\sqrt\Delta$增长；抽样模式还增加了
“哪一种拟合趋势适用”的不确定性。在分量漂移恒定、模式段内固定的解释性情形下，
混合协方差为$u\sum_m\pi_mQ_m+u^2\operatorname{Cov}_{\pi}(b_m)$：
分量内扩散与分量间漂移差异可以按不同速率增长。
HSM的漂移实际依赖位置和太阳条件，因此该式用于解释组成，
不等于实际多步方差，也不证明预测区域具有可靠覆盖。

\subsubsection{训练分组、参数估计与验证选择}\label{sec:training-methods}
\paragraph{HSM：按轨迹段方向硬分组。}
开发段按登记观测间隔$\tau$重采样；不完整的末端间隔只从拟合中去除，评分观测保持原样。
在$\tau=60$个保存秒下，328个训练、76个适应、81个验证窗口分别提供
26,009、5,762、6,149个完整转移。转移目标是二维位移率
$y_i=(x_{i+1}-x_i)/\tau$。
每个拟合批次按段净位移方向排序，再分为三个段数近似平衡的组，同一段的转移继承该组。
这是批次内的方向排序，不是三个固定地理角度区间；训练与适应批次各自重新分组。
这些组标识回归分量，不是经外部验证的步行、转弯或交通状态。
对分量$m$，特征矩阵$D_m$与位移率矩阵$Y_m$给出基础估计
\begin{align}
 \widehat\Theta_m&=\underset{\Theta}{\operatorname{argmin}}
 \left\{\tfrac12\|D_m\Theta-Y_m\|_F^2+
             \tfrac\rho2\|\Theta\|_F^2\right\},
 \qquad \rho=10^{-6},\label{eq:main-full-fit}\\
 \widehat R_m&=\frac1{n_m-1}\sum_{i\in m}
 (e_i-\bar e_m)(e_i-\bar e_m)^\top+10^{-8}I_2,
 \quad e_i=y_i-\widehat\Theta_m^\top\psi_i .
 \label{eq:main-full-residual}
\end{align}
Frobenius范数累加各元素平方，全部系数包括截距均受惩罚。
$n_m$是该组转移数，基础混合权重为$n_m/\sum_k n_k$；
$R_m$描述中心化位移率残差变化，对角小量用于稳定协方差计算，不能校准预测区域。
这些是回归估计，不是连续时间路径似然估计。

\paragraph{RDM：按转移残差方向硬分组并重拟合。}
先在全批次上用同一特征和岭回归拟合$\widehat\Theta_{\rm pool}$，计算
$r_i=y_i-\widehat\Theta_{\rm pool}^{\top}\psi_i$，再一次性指定
\begin{equation}\label{eq:rdm-grouping}
 g_i=\arg\max\{r_{i,e},\ r_{i,n},\ -r_{i,e}-r_{i,n}\}.
\end{equation}
三个得分对应的索引就是三个组；同一轨迹段中的不同转移可以分到不同组。
随后各组仍以原始位移率$y_i$为目标，按式\eqref{eq:main-full-fit}--\eqref{eq:main-full-residual}
分别拟合一套完整的漂移系数和残差协方差，而不是只拟合残差再叠加到共同漂移。
全批次回归仅用于产生标签，不作为共同漂移保留。分组是确定性的硬分组，没有EM迭代或软责任度。
因此，基础拟合确实可以概括为“按规则分三组，各组分别训练一个回归分量”；
两者共享特征设计和正则强度，不共享组内拟合系数。

\paragraph{两者共同的默认训练／适应程序。}
HSM-01与RDM都分别拟合训练与适应批次，再按0.35和0.65混合对应分量的系数、协方差和比例。
这不是合并两批数据后做一次拟合，也不是从训练参数出发的梯度微调。
验证阶段固定漂移，从$s\in\{0.55,0.75,1,1.3,1.7\}$中选择协方差尺度；
保存的$R_m$已经包含$s^2$缩放。
连续时间嵌入使用$Q_m=\tau R_m$：在常系数时长$\tau$下，
$Q_m\tau=R_m\tau^2$与位移率回归隐含的位移协方差一致。
这解释时间因子的来源，不证明它等价于SDE最大似然估计。
参数在评估起点和种子之间保持固定；最终分量概率包含上述训练／适应混合。
HSM保存分量概率为0.3350、0.3302和0.3348，
不是从每个新前缀重新估计的后验。
HSM／RDM标识上述模型构造与预测程序；scratch、Reptile、drift-only、two-step等
另行改变训练／适应策略或参数更新范围，不能与这两个程序名当作同一层级的训练方法清单。
